problem:  
Let \((X,d,\mathfrak{m})\) be a compact geodesic metric measure space satisfying the curvature–dimension condition \(CD(0,N)\) for some finite \(N>1\).  
Assume:  
1. For every \(x\in X\), the measured tangent cone \(T_xX\) (in the pointed measured Gromov–Hausdorff sense) is uniquely Euclidean, i.e. isometric to \((\mathbb{R}^k,|\cdot|,\lambda^k)\) for a fixed integer \(k\) independent of \(x\);  
2. Every Lipschitz function on \(X\) admits a differentiable structure compatible with this Euclidean tangent structure (in Cheeger’s sense).  

Is it true that under these assumptions \(X\) must already satisfy the infinitesimal Hilbertianity condition, and hence be an \(RCD(0,N)\) space, thus isometric (as a metric measure space) to a smooth Riemannian manifold \((M^k,g,\mathrm{vol}_g)\) with nonnegative Ricci curvature?  

Equivalently: do the Euclidean tangent structure and the CD(0,N) synthetic lower Ricci bound together automatically rule out genuinely Finsler behavior?  

Why is it a "good" problem:  
This formulation eliminates ambiguity by fixing dimension and clarifying the metric–measure isomorphism, and it focuses directly on the conceptual boundary between **synthetic curvature bounds** and **analytic structure recovery**.  
It asks whether infinitesimal Hilbertianity—and hence the passage from \(CD\) to \(RCD\)—is compelled by Euclidean tangents alone. This lies precisely at the unknown interface between **Finsler and Riemannian realms** of curvature–dimension theory.  
Resolving it would determine whether tangent-cone rigidity forces global Hilbertian structure, bridging geometric measure theory, synthetic Ricci curvature, and analysis on metric spaces. The problem is simple to state but profoundly deep: it tests whether “being infinitesimally Euclidean everywhere” is enough to recover all features of smooth Ricci geometry. Such a result, positive or negative, would redefine our understanding of the fine topology of curvature bounds in nonsmooth spaces.